
\begin{thebibliography}{99}

\bibitem{lindsey-2026-introspection}
Lindsey, J. (2026). Emergent Introspective Awareness in Large Language Models.
\textit{arXiv:2601.01828v1 [cs.CL]}, 5 Jan 2026. Anthropic.

\bibitem{elhage-2021-circuits}
Elhage, N., Nanda, N., Olsson, C., Henighan, T., Joseph, N., Mann, B., et al. (2021).
A Mathematical Framework for Transformer Circuits. \textit{transformer-circuits.pub}.

\bibitem{yang-2018-softmax}
Yang, Z., Dai, Z., Salakhutdinov, R. and Cohen, W. W. (2018). Breaking the Softmax
Bottleneck: A High-Rank RNN Language Model. \textit{ICLR 2018}. arXiv:1711.03953.

\bibitem{elhage-2022-superposition}
Elhage, N., Hume, T., Olsson, C., Schiefer, N., Henighan, T., Kravec, S., et al. (2022).
Toy Models of Superposition. \textit{arXiv:2209.10652}.

\bibitem{vaswani-2017-attention}
Vaswani, A., Shazeer, N., Parmar, N., Uszkoreit, J., Jones, L., Gomez, A. N.,
Kaiser, {\L}. and Polosukhin, I. (2017). Attention Is All You Need.
\textit{Advances in Neural Information Processing Systems 30}. arXiv:1706.03762.

\bibitem{brown-2020-gpt3}
Brown, T. B., Mann, B., Ryder, N., et al. (2020). Language Models are Few-Shot
Learners. \textit{NeurIPS 33}. arXiv:2005.14165.

\bibitem{kaplan-2020-scaling}
Kaplan, J., McCandlish, S., Henighan, T., et al. (2020). Scaling Laws for Neural
Language Models. \textit{arXiv:2001.08361}.

\bibitem{hoffmann-2022-chinchilla}
Hoffmann, J., Borgeaud, S., Mensch, A., et al. (2022). Training Compute-Optimal
Large Language Models. \textit{arXiv:2203.15556}.

\bibitem{kalai-vempala-2023-calib}
Kalai, A. T. and Vempala, S. S. (2023). Calibrated Language Models Must
Hallucinate. \textit{arXiv:2311.14648}.

\bibitem{kalai-2025-hallucinate}
Kalai, A. T., Nachum, O., Vempala, S. S. and Zhang, E. (2025). Why Language
Models Hallucinate. \textit{arXiv:2509.04664}.

\bibitem{kadavath-2022-know}
Kadavath, S., Conerly, T., Askell, A., et al. (2022). Language Models (Mostly)
Know What They Know. \textit{arXiv:2207.05221}.

\bibitem{cheng-2024-know-gaps}
Cheng, Q., Sun, T., Liu, X., et al. (2024). Can AI Assistants Know What They
Don't Know? \textit{arXiv:2401.13275}.

\bibitem{ferrando-2024-know-entity}
Ferrando, J., Obeso, O., Rajamanoharan, S. and Nanda, N. (2024). Do I Know This
Entity? Knowledge Awareness and Hallucinations in Language Models.
\textit{arXiv:2411.14257}.

\bibitem{zou-2023-repeng}
Zou, A., Phan, L., Chen, S., et al. (2023). Representation Engineering: A
Top-Down Approach to AI Transparency. \textit{arXiv:2310.01405}.

\bibitem{turner-2023-actadd}
Turner, A. M., Thiergart, L., Udell, D., Leech, G., Mini, U. and MacDiarmid, M.
(2023). Activation Addition: Steering Language Models Without Optimization.

\bibitem{chen-2024-selfie}
Chen, H., Vondrick, C. and Mao, C. (2024). SelfIE: Self-Interpretation of Large
Language Model Embeddings. \textit{arXiv:2403.10949}.

\bibitem{ghandeharioun-2024-patchscopes}
Ghandeharioun, A., Caciularu, A., Pearce, A., Dixon, L. and Geva, M. (2024).
Patchscopes: A Unifying Framework for Inspecting Hidden Representations of
Language Models. \textit{arXiv:2401.06102}.

\bibitem{davidson-2024-selfrecog}
Davidson, T. R., Surkov, V., Veselovsky, V., West, R. and Gulcehre, C. (2024).
Self-Recognition in Language Models. \textit{arXiv:2407.06946}.

\bibitem{wang-2025-out-of-context}
Wang, A., Engels, J. and Clive-Griffin, O. (2025). Simple Mechanistic
Explanations for Out-Of-Context Reasoning. \textit{arXiv:2507.08218}.

\bibitem{binder-2024-looking-inward}
Binder, F. J., Chua, J., Korbak, T., et al. (2024). Looking Inward: Language
Models Can Learn About Themselves by Introspection. \textit{arXiv:2410.13787}.

\bibitem{betley-2025-propensities}
Betley, J., Bao, X., Soto, M., Sztyber-Betley, A., Chua, J. and Evans, O.
(2025). Tell Me About Yourself: LLMs Are Aware of Their Learned Behaviors.
\textit{arXiv:2501.11120}.

\bibitem{ji-an-2025-metacog}
Ji-An, L., Xiong, H.-D., Wilson, R. C., Mattar, M. G. and Benna, M. K. (2025).
Language Models Are Capable of Metacognitive Monitoring and Control of Their
Internal Activations. \textit{arXiv:2505.13763}.

\bibitem{plunkett-2025-self-interp}
Plunkett, D., Morris, A., Reddy, K. and Morales, J. (2025). Self-Interpretability:
LLMs Can Describe Complex Internal Processes that Drive Their Decisions, and
Improve with Training. \textit{arXiv:2505.17120}.

\bibitem{comsa-shanahan-2025}
Comsa, I. M. and Shanahan, M. (2025). Does It Make Sense to Speak of
Introspection in Large Language Models? \textit{arXiv:2506.05068}.

\bibitem{song-2025-introspect-language}
Song, S., Hu, J. and Mahowald, K. (2025). Language Models Fail to Introspect
About Their Knowledge of Language. \textit{arXiv:2503.07513}.

\bibitem{song-privileged-2025}
Song, S., Lederman, H., Hu, J. and Mahowald, K. (2025). Privileged Self-Access
Matters for Introspection in AI. \textit{arXiv:2508.14802}.

\bibitem{lindsey-2025-biology}
Lindsey, J., Gurnee, W., Ameisen, E., et al. (2025). On the Biology of a Large
Language Model. \textit{Transformer Circuits Thread}, Anthropic.

\bibitem{kammerer-frankish-2023}
Kammerer, F. and Frankish, K. (2023). What Forms Could Introspective Systems
Take? A Research Programme. \textit{Journal of Consciousness Studies} 30(9--10):13--48.

\bibitem{long-2023-introspection}
Long, R. (2023). Introspective Capabilities in Large Language Models.
\textit{Journal of Consciousness Studies} 30(9--10):143--153.

\bibitem{laine-2024-sad}
Laine, R., Chughtai, B., Betley, J., et al. (2024). Me, Myself, and AI: The
Situational Awareness Dataset (SAD) for LLMs. \textit{Advances in Neural
Information Processing Systems} 37:64010--64118.

\bibitem{kamath-2025-concordance}
Kamath, H., Ameisen, E., Kauvar, I., et al. (2025). Tracing Attention
Computation Through Feature Interactions. \textit{Transformer Circuits Thread},
Anthropic.

\bibitem{chalmers-2023-conscious}
Chalmers, D. J. (2023). Could a Large Language Model be Conscious?
\textit{arXiv:2303.07103}.

\bibitem{butlin-2023-conscious}
Butlin, P., Long, R., Elmoznino, E., et al. (2023). Consciousness in Artificial
Intelligence: Insights from the Science of Consciousness.
\textit{arXiv:2308.08708}.

\bibitem{albantakis-2023-iit}
Albantakis, L., Barbosa, L., Findlay, G., et al. (2023). Integrated Information
Theory (IIT) 4.0: Formulating the Properties of Phenomenal Existence in
Physical Terms. \textit{PLoS Computational Biology} 19(10):e1011465.

\bibitem{baars-1993}
Baars, B. J. (1993). \textit{A Cognitive Theory of Consciousness}. Cambridge
University Press.

\bibitem{baars-1997}
Baars, B. J. (1997). In the Theatre of Consciousness: Global Workspace Theory,
a Rigorous Scientific Theory of Consciousness. \textit{Journal of Consciousness
Studies} 4(4):292--309.

\bibitem{rosendahl-2005}
Rosenthal, D. M. (2005). \textit{Consciousness and Mind}. Clarendon Press.

\bibitem{block-1995}
Block, N. (1995). On a Confusion About a Function of Consciousness.
\textit{Behavioral and Brain Sciences} 18(2):227--247.

\bibitem{tononi-2004}
Tononi, G. (2004). An Information Integration Theory of Consciousness.
\textit{BMC Neuroscience} 5:42.

\bibitem{lau-rosenthal-2011}
Lau, H. and Rosenthal, D. (2011). Empirical Support for Higher-Order Theories
of Conscious Awareness. \textit{Trends in Cognitive Sciences} 15(8):365--373.

\bibitem{carruthers-2017}
Carruthers, P. (2017). Higher-Order Theories of Consciousness. In
\textit{The Blackwell Companion to Consciousness}, pp.~288--297.

\bibitem{seth-2024}
Seth, A. K. (2024). Conscious Artificial Intelligence and Biological
Naturalism. \textit{Behavioral and Brain Sciences}, pp.~1--42.

\bibitem{searle-1992}
Searle, J. R. (1992). \textit{The Rediscovery of the Mind}. MIT Press.

\bibitem{long-2024-welfare}
Long, R., Sebo, J., Butlin, P., et al. (2024). Taking AI Welfare Seriously.
\textit{arXiv:2411.00986}.

\bibitem{shannon-1948}
Shannon, C. E. (1948). A Mathematical Theory of Communication.
\textit{Bell System Technical Journal} 27:379--423, 623--656.

\bibitem{cover-thomas-2006}
Cover, T. M. and Thomas, J. A. (2006). \textit{Elements of Information Theory},
2nd edition. Wiley.

\bibitem{kolmogorov-1965}
Kolmogorov, A. N. (1965). Three Approaches to the Quantitative Definition of
Information. \textit{Problems of Information Transmission} 1(1):1--7.

\bibitem{valiant-1984}
Valiant, L. G. (1984). A Theory of the Learnable. \textit{Communications of the
ACM} 27(11):1134--1142.

\bibitem{wolpert-macready-1997}
Wolpert, D. H. and Macready, W. G. (1997). No Free Lunch Theorems for
Optimization. \textit{IEEE Transactions on Evolutionary Computation}
1(1):67--82.

\bibitem{shalev-shwartz-2014}
Shalev-Shwartz, S. and Ben-David, S. (2014). \textit{Understanding Machine
Learning: From Theory to Algorithms}. Cambridge University Press.

\bibitem{mcallester-1999}
McAllester, D. A. (1999). Some PAC-Bayesian Theorems. \textit{Proceedings of
the Eleventh Annual Conference on Computational Learning Theory (COLT)}.

\bibitem{rice-1953}
Rice, H. G. (1953). Classes of Recursively Enumerable Sets and Their Decision
Problems. \textit{Transactions of the American Mathematical Society}
74(2):358--366.

\bibitem{gold-1967}
Gold, E. M. (1967). Language Identification in the Limit.
\textit{Information and Control} 10(5):447--474.

\bibitem{pearl-2009}
Pearl, J. (2009). \textit{Causality: Models, Reasoning, and Inference}, 2nd
edition. Cambridge University Press.

\bibitem{friston-2010}
Friston, K. (2010). The Free-Energy Principle: A Unified Brain Theory?
\textit{Nature Reviews Neuroscience} 11(2):127--138.

\bibitem{dayan-abbott-2001}
Dayan, P. and Abbott, L. F. (2001). \textit{Theoretical Neuroscience:
Computational and Mathematical Modeling of Neural Systems}. MIT Press.

\bibitem{manheim-garrabrant-2018-goodhart}
Manheim, D. and Garrabrant, S. (2018). Categorizing Variants of Goodhart's Law.
\textit{arXiv:1803.04585}.

\bibitem{skalse-2022-reward-hacking}
Skalse, J., Howe, N., Krasheninnikov, D. and Krueger, D. (2022). Defining and
Characterizing Reward Hacking. \textit{arXiv:2209.13085}.

\bibitem{joudaki-2025-plasticity}
Joudaki, A., Lanzillotta, G., Razlighi, M. S. and Mirzadeh, I. (2025). Barriers
for Learning in an Evolving World: Mathematical Understanding of Loss of
Plasticity. \textit{arXiv:2510.00304}.

\end{thebibliography}
